

\documentclass[12pt]{article}

\usepackage[margin = .8in]{geometry}
\usepackage{amsmath}
\usepackage{graphicx}
\usepackage{multicol, enumerate, tabularx}

\usepackage[parfill]{parskip}

\usepackage{adjustbox, soul}

\usepackage{fancyhdr}
\pagestyle{fancy}

\lhead{\emph{Math F113X: Double Transposition}}
\rhead{\emph{lecture notes}}

\renewcommand{\emph}[1]{\textsf{\textbf{#1}}}
\usepackage{tikz}
\usetikzlibrary{calc,trees,positioning,arrows,fit,shapes,through, backgrounds}
\usetikzlibrary{patterns}

\usetikzlibrary{decorations.markings}
\usetikzlibrary{arrows}

\usepackage{pgfplots}

\usepackage{longtable}
\usepackage{tabularx}

\newcommand{\ds}{\displaystyle}
\newcommand{\ans}[1][1in]{\rule{#1}{.5pt}}

\newcommand{\points}[1]{(#1 points.)}		% Trying to be lazy.

\usepackage{array}
\newcolumntype{L}[1]{>{\raggedright\let\newline\\\arraybackslash\hspace{0pt}}m{#1}}
\newcolumntype{C}[1]{>{\centering\let\newline\\\arraybackslash\hspace{0pt}}m{#1}}
\newcolumntype{R}[1]{>{\raggedleft\let\newline\\\arraybackslash\hspace{0pt}}m{#1}}
\newcommand{\red}[1]{\textcolor{red}{#1}}

\newcommand{\be}{\begin{enumerate}}
\newcommand{\ee}{\end{enumerate}}

%\topmargin -1in
%\textheight 9.5in
%\oddsidemargin -0.3in
%\evensidemargin \oddsidemargin
%\pagestyle{empty}
%%\marginparwidth 0.5in
%\textwidth 7in
%\parindent 0in

%--------------------------------------------------------------------------------------------------------------------------------------------------------------------------
%						Document
%--------------------------------------------------------------------------------------------------------------------------------------------------------------------------


\begin{document}
%\pagestyle{fancy}
%\begin{center}
%{\Large  Cryptography (Day 3)}
%\end{center}
\begin{enumerate}
\item \emph{Double Transposition Cipher:} Use two keywords (often of different lengths). Do a transposition cipher with the first keyword to produce a first ciphertext. If there are spaces, ignore them. Then encrypt that ciphertext (as though it were plaintext) using the second keyword.

\emph{Encryption}

\be[i.]
\item Write the plaintext in rows.
\item Permute the columns using the keyword.
\item Read off the \emph{columns} to form the new ``plaintext'' 
\item Write it in rows in the second grid.
\item Permute the columns using the second keyword.
\item Read off the \emph{columns} to form the final ciphertext.
\ee


\emph{Example:} Encrypt the word TRANSPOSITION using the first key {\tt KEY} and the second key {\tt SHORT}.


\newpage


\emph{Decryption:}

\be[i.]
\item Count the number of characters. Determine how many ``extra'' characters there will be for the second keyword. 
\item Write your second grid with the permuted keyword. 
\item Fill in the \emph{columns} of the grid, making sure to put the extra characters in the appropriate columns. Mark those columns with *. For example, if your second keyword is {\tt SHORT} and you have two ``extra'' characters, then columns S and H get one more letter than the other columns.
\item (Un)permute the columns.
\item Read off the {\bf rows}.
\item Using the number of characters, determine the number of ``extra'' characters for the first keyword.
\item Fill in the columns of the grid, filling in the extra long columns (marked with *) using the sorted keyword. For example, if your first keyword is {\tt KEY} and you have two ``extra'' characters, then columns K and E get one more letter than the other columns.
\item Unpermute the columns.
\item Read off the {\bf rows}
\ee


\emph{Decrypt} what we got when we encrypted TRANSPOSITION.


\vfill


\emph{Decrypt} the ciphertext \so{ETETS PIHXC R} assuming it was encrypted using the scheme above.

\vfill



\end{enumerate}
\end{document}

%-------------------------------------------------------------------------------------------------------------------------------------------------------------------------------------------------------------------

%%% Local Variables:
%%% mode: latex
%%% TeX-master: t
%%% End:
